\section{讨论}
\label{sec:discussion}

扩大病例规模能换来多少精度，取决于估计目标。瞬时再生数的规模收益不饱和，方差随规模持续下降（命题~\ref{prop:current}）。跨期平均再生数的误差含两类不随规模下降的方差：环境方差 $v_R$ 来自传播条件的逐期波动，报告异质性方差 $v_H$ 来自报告比例的逐期波动（命题~\ref{prop:average} 与命题~\ref{prop:report-noise}）。方差平衡规模给出规模收益开始明显减小的病例数，报告率波动时它由 $A/v_R$ 降为 $A_H/(v_R+v_H)$。Parag 等\citep{parag2022quantifying, parag2022fundamental}的信息论分析刻画了随规模下降的那一部分，在离散参数趋于无穷时与命题~\ref{prop:current} 一致。生态学中环境随机性与人口随机性的区分\citep{engen1998demographic, rohani2002determinism}由此在监测设计中有了一个以病例数计的对应量。报告异质性则是监测系统自身带来的第二个不可约成分，仅凭计数的前两阶矩无法把它与环境方差分开。

本文的估计目标与 Eales 和 Riley（\citeyear{eales2024differences}）讨论的表观再生数不同。后者由一条病例序列和给定世代间隔反推，是真实再生数的滞后加权平均；本文的估计目标1是条件瞬时值，估计目标2是跨期均值。使用滞后计数构造增长统计量时，两者都会受到卷积影响，这是本文把方法放在最后、并在更新过程一节单独处理它的原因。反卷积可以减轻表观再生数的滞后，但不能消除估计目标2中的环境方差；后者不是估计方法的误差，而是目标本身的波动。

上述下确界不约束允许条件偏差的估计量。收缩估计、经验贝叶斯估计\citep{efron1975data}以及 EpiFilter 一类借用相邻时期信息的平滑方法\citep{parag2021improved}，都以对 $R_t$ 的条件偏差换取方差。以补充材料 S1.11 节的线性收缩为例，取图~\ref{fig:variance} 的参数与 $m=50$，$V=0.0825$；先验中心偏离 $\Rbar$ 不超过 0.279 时，最优收缩的均方误差就低于 $v_R=0.04$，偏离 0.1 时为 0.0089，约为 $v_R$ 的 22\%。代价是估计量向先验中心偏移，偏差随当期环境变化，在 $R_t$ 远离常态的时期最大，而这正是预警最需要准确的时候。本文关于下确界与方差平衡规模的结论只覆盖条件无偏的单期估计和 H4 下的一步预测；平滑估计的均方误差如何随规模变化，需要另行推导。

在监测设计上，三类杠杆对应不同的投入。扩大单期规模降低抽样项，直到误差接近由 $v_R+v_H$ 决定的单期误差下限；它对预警始终有价值，对趋势估计的价值在规模达到方差平衡规模的量级后明显减小。增加观测期数或汇总更多辖区降低环境项；规模足够大以后，$v_R$ 本身的估计精度也只随期数提高（命题~\ref{prop:vR-info}）。统一报告口径，或在估计中使用观测到的报告率，去除 $v_H$。在命题~\ref{prop:levers} 的例子中，期数增至四倍把均方误差从 22.5 降到 6.2（单位 $10^{-3}$），扩大单期规模十倍只降到 18.7（图~\ref{fig:r3-levers}C--D）。报告偏差在各期持续时，增加期数对报告项无效，只有统一口径能处理。期数与规模之间的取舍还取决于环境的时间相关：$R_t$ 正相关时，预算低于阈值 $\mathcal B^\dagger$ 应设更少、更大的期，高于阈值才应设更多、更小的期（命题~\ref{prop:ar-design}）；按独立情形的近似设计小预算监测，会选出过小的单期规模。

对监测实践，预警和趋势估计宜采用不同的设计。判断本周疫情是否在增长属于估计目标1，规模的价值持续存在，小规模时应使用精确检验，并同时报告首次达到和稳定保持目标功效所需的规模（第~\ref{sec:res-test}~节），资源应投向报告及时、覆盖完整的单元。估计平均传播水平属于估计目标2，规模达到方差平衡规模量级后，增加观测期数或汇总多辖区更有效（式~\eqref{eq:ar-sum} 与式~\eqref{eq:pooling}），全国加总序列结果与空间推论一致（第~\ref{sec:res-seasons}~节）。报告率固定且很低时，提高报告率与扩大覆盖的效果相差无几（注~\ref{rem:reporting}），取舍主要取决于成本；报告率随期波动时，二项型波动在规模足够大时使提高报告率占优，制度型波动下扩大覆盖仍占优，两条途径都不改变 $v_H$（推论~\ref{cor:switch}）。NHSN 的填报医院比例在流行初期波动很小，自愿填报期则大了两个数量级（第~\ref{sec:res-reporting-adj}~节），报告标准化的收益因此主要出现在填报制度变化的时期；切换制度时应同时记录各单位的填报状态，以便事后按报告率调整。

流行阶段转换时出现了另一种误差来源。达峰期以初期均值为基准时，绝大部分损失来自基准偏移，扩大规模无法消除它；换用峰期自身的基准后，规模项重新为正（第~\ref{sec:res-seasons}~节）。阶段转换时应优先更新基准模型，同时继续维持病例规模。FluSight 集成预测在峰期的偏差占比较低，提示能跟踪阶段变化的预测模型可以压低这一项；该对比来自不同的模型和样本，只有方向上的参考价值。报告制度切换的证据需要按流行季解读。全样本上第二期稳健离散度不变的表象，来自第二期季内州周所占比例很低；分层后，季外三期几乎相同，季内第二期的离散度明显扩大。季内证据全部来自自愿填报期的最后一个月，区间很宽，填报不全造成的跳变与该月本身的流行特征无法分开。数据质量监控应按流行季分别进行。

在实证解读上，主要局限在于增长统计量的口径。流感的世代间隔短于一周，$G_t$ 是两代左右增长的合成，其均值与方差都不同于单代模型中的 $R_t$。实证分析因此只检验了误差随规模变化的形状，拟合的 $a$ 和 $b$ 不能解读为 $v_R/\Rbar^2$ 与 $A/\Rbar^2$ 的估计，经验比值 $b/a$ 与情景区间的吻合也只在量级上成立；估计单代参数需要日度数据或显式的世代间隔反卷积。均值形状比较支持误差含不随规模消失的部分，但在现有样本覆盖下，严格底板与指数约 0.27 的幂律衰减拟合相近，需要更宽规模才能区分两者。截距与斜率的绝对水平也都只得到部分解释：斜率比只表明分母含有报告计数噪声，斜率的绝对水平约有一半能由二项分母噪声解释，截距比可以与平稳环境模型区分，却与增长期、含入院延迟的更新过程以及乘性报告噪声的范围相交。跨病原分级为复现 0、部分复现 4、方向一致 3，证据主要来自截距超出量，不能视为独立复制。

流行季的数目是最根本的限制。按推论~\ref{cor:season}，季层面方差的估计精度由季数决定，增加州数、周数或报告规模都不改变它。实证结果与此一致：截距的重抽标准差在规模外推到无穷大时几乎不变，把季作为随机效应时区间很宽。按州聚类的区间因此只描述给定已观测三个季时的不确定性。对截距水平的进一步认识需要积累更多流行季，扩大单季样本的作用有限。跨病原与样本外季的复制都来自同一个 NHSN 文件，共用同一套报告系统，系统层面的报告特征无法由这些复制排除。方差分解还依赖 H5，即报告率 $p_t$ 与再生数 $R_t$ 相互独立；若检测或填报行为随流行强度变化，报告波动会与 $R_t$ 相关，协方差项会混入截距。本文的报告调整只覆盖可观测的填报医院比例，不覆盖随流行强度变化的就诊、检测与收治行为，可观测报告成分贡献很小的结论不能推广到报告行为整体。工具变量回归的排除限制无法检验；流行初期窗口只有八到十一周，按州聚类的重抽在聚类数较少时方差估计较不稳定\citep{cameron2011robust, cameron2015practitioner}。情景区间依赖 $k$ 与住院比例的外部取值，精确增长检验要求 $k$ 已知；缺失州周按零处理，而未填报与零住院并不等同。这些限制不改变方差分解的方向，但限定了具体数值的可移植性。
